\chapter{Related Work}
\label{ch:related-work}

This chapter reviews four strands of prior work and then synthesises the gap this thesis fills.
Section~\ref{sec:rw-methods} surveys machine unlearning and its method families;
Section~\ref{sec:rw-benchmarks} surveys the benchmarks and metrics used to evaluate it, and observes that
they are almost entirely output-based; Section~\ref{sec:rw-robustness} collects the mounting evidence that
these output metrics are fragile because unlearning tends to \emph{suppress} rather than \emph{erase};
Section~\ref{sec:rw-mechinterp} reviews mechanistic interpretability and the lens techniques; and
Section~\ref{sec:rw-gap} shows that interpretability has met unlearning only on the \emph{intervention} side,
leaving the \emph{evaluation} side open, which is precisely the gap this thesis addresses.

\section{Machine Unlearning and its Method Families}
\label{sec:rw-methods}

The problem of removing the influence of specific training data from a trained model was formalised for
classical models as exact unlearning, most influentially through the SISA framework, which shards and
slices the data so that only the affected shard need be retrained \citep{bourtoule2021machine,cao2015forget}.
Exact retraining is intractable for LLMs, so the literature has shifted to \emph{approximate} unlearning via
fine-tuning; recent surveys map the resulting landscape of formulations, methods, metrics, and open problems
\citep{liu2025rethinking}.

Within this landscape, four families recur. \emph{Gradient-based} methods maximise the loss on the forget
set, either directly (gradient ascent, \citealp{jang2022knowledge,yao2023llmunlearning}) or in combination
with retain-set descent (Gradient Difference); they are simple but unstable. \emph{Preference-based} methods,
of which Negative Preference Optimization is the leading example, cast forgetting as preference optimisation
over dispreferred completions and enjoy a bounded, collapse-resistant objective \citep{zhang2024npo,
rafailov2023dpo}. \emph{Refusal / redirection} methods, such as the ``I don't know'' baseline, change the
output policy on the forget distribution rather than the underlying representation \citep{maini2024tofu}, a
strategy also embodied by the ``Who's Harry Potter'' approach of replacing target content with generic
alternatives \citep{eldan2023whp}. \emph{Representation-based} methods intervene on internal activations
directly; Representation Misdirection for Unlearning steers forget-set hidden states towards noise at chosen
layers while preserving retain-set activations \citep{li2024wmdp}. This thesis studies one representative of
each family (GradDiff, NPO, IDK, RMU), because a central claim, that the four leave distinct \emph{internal}
signatures despite looking similar at the output, only makes sense across the whole taxonomy.

\section{Benchmarks and the Evaluation of Unlearning}
\label{sec:rw-benchmarks}

Progress in unlearning is measured on a handful of benchmarks, and it is instructive that they agree on
\emph{what} to measure. The model's output. TOFU evaluates unlearning of fictitious author profiles via
``forget quality'', a distributional test on answer probabilities, together with model utility, and reports
that none of its baselines achieve effective unlearning \citep{maini2024tofu}. MUSE proposes a six-way
output evaluation spanning verbatim and knowledge memorisation, privacy leakage, and utility
\citep{shi2024muse}. WMDP measures hazardous knowledge through multiple-choice accuracy and introduces RMU
as a mitigation \citep{li2024wmdp}. RWKU targets real-world entity knowledge with membership-inference and
adversarial probes \citep{jin2024rwku}. Across all of them the metrics are lexical or behavioural, such as
exact-match, ROUGE, accuracy, gold-token probability, and membership-inference scores.

The most relevant recent effort is OpenUnlearning, which unifies thirteen methods and sixteen metrics across
TOFU, MUSE, and WMDP and, crucially, proposes a \emph{meta-evaluation}, a way to score the metrics
themselves on their \emph{faithfulness} (a faithful metric assigns higher scores to models that hold the
targeted knowledge than to those that lack it) and \emph{robustness} (the score should not improve under
benign changes such as quantisation, nor faster than a retain model under adversarial ones)
\citep{dorna2025openunlearning}. OpenUnlearning thereby supplies the yardstick this thesis adopts, but it
applies that yardstick only to existing output and lexical metrics; it does not introduce a metric that
reads the model's internal state. Independent work on semantic unlearning reaches a compatible conclusion
from the method side, arguing that lexical-level forgetting leaves knowledge recoverable by paraphrase and
proposing a semantic-residual measure of leakage as a more faithful evaluation than traditional metrics
\citep{semantic2025unlearning}.

Hubble is the benchmark on which this thesis builds \citep{allegro2025hubble}. Its controlled insertion of
facts yields, for each fact, a never-knew reference and a knows reference, converting unlearning evaluation
from guesswork into a measurement against ground truth. Two design choices are directly pertinent here.
First, Hubble deliberately \emph{untied} its output embeddings to support the logit and tuned lenses,
citing exactly the interpretability literature this thesis employs. Second, its concluding research
questions explicitly invite interpretability studies that probe whether inserted knowledge is localised, and
call for more intuitive, robust, and \emph{validated} memorisation metrics. Hubble thus provides both the
substrate and the stated motivation for an internal metric, while itself evaluating only at the output.

\section{Does Unlearning Really Erase? Robustness and Attacks}
\label{sec:rw-robustness}

A distinct and rapidly growing literature stress-tests unlearned models and repeatedly finds that the
``forgotten'' knowledge is still present, merely suppressed. Quantising an unlearned model back to lower
precision recovers a large fraction of the target knowledge, and the authors explicitly attribute the
failure of existing benchmarks to their output-only view. Current evaluations ``fail to detect'' whether a
method truly forgets or merely hides \citep{zhang2024quantization}. Relearning attacks that fine-tune on a
small amount of related data, or even supply it in context, revive the target, motivating unlearning
objectives designed to be resilient to such relearning \citep{lynch2024eight,fan2025relearn}. Black-box
evaluations and weight-space probes reach the same verdict, recovering supposedly removed information and
questioning whether current methods remove it from the weights at all \citep{thaker2024blackbox,
deeb2024weights}. Lynch et al.\ systematise this into eight complementary robustness checks, including
other-language queries, jailbreaks, in-context and fine-tuning relearning, and, notably, direct probes of
\emph{latent knowledge} \citep{lynch2024eight}.

The single most relevant result for this thesis probes latent knowledge with an interpretability tool.
Patil et al.\ show that even after state-of-the-art model \emph{editing} such as ROME, traces of the
``deleted'' fact persist in intermediate hidden states, and can be recovered by whitebox logit-lens attacks
that read out the top tokens at each layer; ``deleted'' information is recovered a substantial fraction of
the time \citep{patil2023candelete}. This establishes, directly and internally, that suppression at the
output can coexist with an intact internal representation. However, that work targets knowledge \emph{editing}
rather than the modern unlearning methods, treats the logit lens as an adversarial \emph{attack} rather than
a systematic readout across a controlled benchmark, and does not distil its observations into a validated
\emph{metric} of unlearning. Those three limitations mark the space this thesis occupies.

\section{Mechanistic Interpretability and the Lens Techniques}
\label{sec:rw-mechinterp}

Mechanistic interpretability seeks to explain a network's behaviour in terms of its internal computation.
The residual-stream view of transformers, in which each block reads from and writes to a shared running
representation, provides the substrate for reading intermediate states \citep{elhage2021framework}, and a
line of work has localised factual knowledge to feed-forward key-value memories and identifiable
mid-network sites \citep{geva2021kv,meng2022rome}. To decode an intermediate state into vocabulary space one
applies a \emph{lens}. The \emph{logit lens} applies the model's own unembedding to a hidden state and is
the simplest, parameter-free option \citep{nostalgebraist2020logitlens}. The \emph{tuned lens} fits a small
per-layer affine map so the readout reproduces the final distribution, improving reliability but, because it
is trained to match the output, tending to make the answer appear decodable earlier than the model actually
computes it \citep{belrose2023tunedlens}. The most recent, the \emph{Jacobian lens}, transports a
residual-stream vector at any layer through the model's own averaged input-output Jacobian over a text
corpus before decoding it with the unembedding, yielding an \emph{interventionally-grounded} readout of what
a representation is disposed to make the model say \citep{anthropic2026workspace}. These three form a ladder
of increasing fidelity to the model's own computation, from correlational to calibrated to
interventionally-grounded, and the differences among them become decisive when a lens is used to define a
metric rather than to produce a picture.

\section{Interpretability meets Unlearning, and the Gap}
\label{sec:rw-gap}

Mechanistic interpretability and unlearning have begun to intersect, but so far almost exclusively on the
\emph{intervention} side. Mechanistic Unlearning uses interpretability to \emph{localise} the components
responsible for a fact and edits or unlearns them, finding that localising to genuine factual-recall
mechanisms yields more robust unlearning than localising by output preservation alone
\citep{guo2024mechanistic}; sparse-autoencoder clamping similarly uses interpretable features to perform the
forgetting \citep{sae2025clamping}. In these works interpretability is a means of \emph{doing} unlearning
more robustly, not of \emph{measuring} whether unlearning succeeded. On the evaluation side, the logit lens
has appeared only as an attack on edited models \citep{patil2023candelete}, and the meta-evaluation of
metric faithfulness has been applied only to output and lexical metrics \citep{dorna2025openunlearning}.

Bringing these threads together yields the gap this thesis addresses. First, unlearning is evaluated almost
entirely with output and lexical metrics \citep{maini2024tofu,shi2024muse,li2024wmdp,dorna2025openunlearning}.
Second, a large and consistent body of evidence, from quantisation and relearning attacks to internal
logit-lens probes, shows that these metrics are fooled because unlearning suppresses rather than erases
\citep{zhang2024quantization,lynch2024eight,fan2025relearn,thaker2024blackbox,deeb2024weights,
patil2023candelete,semantic2025unlearning}. Third, although interpretability can \emph{improve} unlearning
\citep{guo2024mechanistic,sae2025clamping}, no prior work turns an interpretability readout into a
\emph{faithful, validated evaluation metric}. Fourth, the tools and the invitation already exist. The tuned
lens is reliable enough to serve as a readout \citep{belrose2023tunedlens}, the Jacobian lens offers an
interventionally-grounded alternative \citep{anthropic2026workspace}, and Hubble was built with untied
embeddings expressly to support
the logit and tuned lenses and calls for validated internal metrics \citep{allegro2025hubble}.

This thesis fills that gap. We integrate a corrected logit lens into Hubble and derive from its rank-by-layer
trajectory a dual-anchor distance, \(D_\text{mech}\), that measures how much a fact's internal trajectory
resembles the ``knows'' rather than the ``never-knew'' reference. We then subject \(D_\text{mech}\) to the
faithfulness meta-evaluation based on OpenUnlearning and show that, on Hubble-8B, it is the only metric,
mechanistic or lexical, that is simultaneously strongly knowledge-discriminative and substantially
paraphrase-stable, a result we then replicate across a second benchmark (OpenUnlearning TOFU) and three
scales. Finally, we characterise all three lenses on the same facts, explaining why the calibrated tuned
lens cannot ground the metric and confirming the taxonomy with the interventionally-grounded Jacobian lens.
In doing so we answer, on the evaluation
side, the question that both the robustness literature and the Hubble benchmark leave open.
